\section{Matched Development and HTTP Follow-up}
The follow-up evaluates the pinned base and existing adapters on the same
1,495 development rows under two interfaces: original training completion
prompts (legacy), and the common chat task specification. All complete paired
runs use bfloat16, batch size eight, greedy decoding, a 64-token output cap and
the strict scorer. Neither training nor final-test scoring occurs. An initial
batch-one run was interrupted for runtime reasons; its partial journal is
retained but excluded. Batching changes are shared by all paired candidates.

\begin{table}[t]
\centering\small
\caption{Matched development primary scores: macro-F1 for relevance and identity;
normalized exact joint accuracy for extraction. Identical rows within each task.}
\label{tab:followup}
\begin{tabular}{llrr}\toprule
Task & Interface & Base & Adapter\\\midrule
Relevance & legacy & 0.0000 & 0.3333 \\
Relevance & chat & 0.2568 & 0.2665 \\
Identity & legacy & 0.0000 & 0.8617 \\
Identity & chat & 0.3110 & 0.8886 \\
Extraction & legacy & 0.0000 & 0.9150 \\
Extraction & chat & 0.5825 & 0.9025 \\
\bottomrule\end{tabular}\end{table}

Interface and adaptation must be interpreted jointly (Table~\ref{tab:followup}).
Strict schema failures remain in the denominator; per-class and validity
results accompany the raw outputs. The recorded historical training-source
audit finds no exact-text matches on these development rows, but identity
entity exposure remains unresolved. Thus this is a retrospective development
comparison, not certified independent generalization. Paired binary/joint
correctness diagnostics use connected groups and 2,000 exploratory bootstrap
resamples; their intervals do not describe macro-F1 effects. No repeated-seed
training effect or matched frontier-API ranking is established.

The unadapted base fails the strict output contract on every legacy-prompt
row. Its zero scores measure interface failure, not absence of commerce knowledge;
the chat comparison is essential to interpretation. Under chat, relevance
macro-F1 changes only from 0.2568 to 0.2665, whereas identity changes from 0.3110
to 0.8886. The adapter predicts no complement labels in either interface.
Extraction joint accuracy changes from 0.5825 to 0.9025 under chat.
The connected grouping yields 33 relevance, 5 identity and 367 extraction
components. We suppress the identity bootstrap interval because fewer than 20 components
do not support that uncertainty estimate. BF16 greedy outputs can depend on
batching; the batch-eight base is rerun rather than substituted by the earlier
batch-one pilot. These findings distinguish task adaptation, label compliance
and interface sensitivity rather than supporting a universal gain.

The NuExtract-2.0-2B comparison uses the native schema interface for text, with
verbatim-string brand/color fields and no demonstrations or prompt search.
On the same 400 extraction rows it obtains normalized joint accuracy
0.4675, normalized brand/color F1
0.5667/0.8283, and invalid rate
0.0000. The revision, prompt and raw outputs are
retained. This closes one missing comparator run, not broader baseline coverage.

\begin{table}[t]
\centering\small
\caption{GPU loopback HTTP prototype. B: maximum batch size; C: client
concurrency. QPS counts completed requests; correctness is over all attempts.
p95 is the mean of three within-run percentiles, not a pooled percentile.}
\label{tab:httpfollowup}
\setlength{\tabcolsep}{3pt}
\begin{tabular}{rrrrr}\toprule
B & C & QPS & Correct & p95 (s)\\\midrule
1 & 1 & 9.01 & 0.900 & 0.116 \\
1 & 4 & 8.96 & 0.900 & 0.466 \\
1 & 8 & 9.06 & 0.900 & 0.933 \\
8 & 1 & 8.57 & 0.900 & 0.124 \\
8 & 4 & 29.11 & 0.905 & 0.148 \\
8 & 8 & 40.56 & 0.900 & 0.287 \\
\bottomrule\end{tabular}\end{table}

The shared identity adapter is served through a bounded HTTP queue on one
RTX 3090, comparing batch limits one and eight with at most 5\,ms batching wait.
Each cell has three repetitions of 200 requests on the same development trace;
50 warmup requests per engine/repetition are excluded. Generation permits eight
tokens. Client and server share the pod, and the load is closed-loop, so these
results do not identify open-loop overload behavior or external network latency.
The smaller workload is an explicit budgeted departure from the prospective
larger serving study; tail estimates are exploratory.

Quality and throughput are reported together (Table~\ref{tab:httpfollowup});
all attempted responses, errors and timestamps are retained. Window-specific
cost estimates use the quoted GPU rate of \$0.22/hour and exclude provisioning,
idle time, storage and network. They are not invoices, production SLO costs or
an API cost advantage. The reference queue is not a new serving algorithm or
a comparison with an established optimized inference engine.

